% main_cn.tex - 中文论文（计算机学报格式）
% Censor: 一种仿生双通路微表情识别框架

\documentclass[12pt,a4paper]{article}

% 中文支持
\usepackage{ctex}
\usepackage{CJKutf8}

% 常用包
\usepackage{amsmath,amssymb}
\usepackage{graphicx}
\usepackage{booktabs}
\usepackage{multirow}
\usepackage{algorithm}
\usepackage{algorithmic}
\usepackage{hyperref}

% 页面设置
\usepackage[top=2.5cm,bottom=2.5cm,left=2.5cm,right=2.5cm]{geometry}

% 参考文献样式
\usepackage[numbers,sort&compress]{natbib}
\bibliographystyle{gbt7714-numerical}

\begin{document}

\title{Censor：一种仿生双通路微表情识别框架}
\author{
作者甲$^{1)$\dagger}$ \quad 作者乙$^{1)}$ \quad 作者丙$^{1)}$ \\
$^{1)}$某某大学计算机科学与技术学院，某某省某某市 000000 \\
$\dagger$通信作者, E-mail: author@example.com
}

\maketitle

% 中文摘要
\begin{abstract}
微表情识别（MER）由于面部运动的细微空间幅度和短暂时持续而面临巨大挑战。本文提出\textbf{Censor}，一种受神经科学梭状回-杏仁核回路启发的仿生双通路神经网络架构。框架包含：快通路（3D ResNet-18处理光流，类比皮层下路径）和慢通路（3D Swin Transformer处理RGB+rPPG，类比皮层路径），通过杏仁核启发的注意力门控和混合专家分类进行融合。在CASME II四类LOSO评估中，Censor达到0.85准确率和0.80 F1分数。跨数据集实验表明从CASME II迁移到SMIC和SAMM可获得0.74--0.76准确率，验证了强泛化能力。消融实验证实了各架构组件的贡献。

\keywords{微表情识别；双通路神经网络；仿生计算；3D Swin Transformer；混合专家；留一受试者交叉验证}
\end{abstract}

% 英文摘要
\begin{enabstract}
Micro-expression recognition (MER) remains challenging due to subtle spatial magnitude and brief temporal duration of involuntary facial movements. This paper presents \textbf{Censor}, a biomimetic dual-pathway neural architecture inspired by the fusiform-amygdala circuit. The framework comprises a fast pathway (3D ResNet-18 on optical flow, analogous to subcortical route) and a slow pathway (3D Swin Transformer on RGB+rPPG, analogous to cortical route), integrated through amygdala-inspired gating and mixture-of-experts classification. On CASME II 4-class LOSO, Censor achieves 0.85 accuracy and 0.80 F1-score. Cross-dataset experiments show 0.74--0.76 accuracy when transferring to SMIC and SAMM. Ablation studies validate each component's contribution.

\enkeywords{Micro-expression recognition; Dual-pathway neural network; Biomimetic computing; 3D Swin Transformer; Mixture of experts; Leave-one-subject-out}
\end{enabstract}

% 正文
\newpage
% introduction_cn.tex - 引言

\section{引言}

面部微表情是个体试图隐藏或抑制真实情感时发生的非自愿、短暂面部运动\cite{ekman1969nonverbal}。与持续0.5--4秒的宏表情不同，微表情仅持续1/25至1/5秒\cite{ekman2003darwin}，具有低强度、部分面部参与和快速起始-顶点-偏移动态特性，即使对训练有素的编码者也难以检测和识别\cite{frank2009training}。

微表情识别的计算建模在过去十年取得显著进展，得益于CASME II\cite{yan2014casme2}、SAMM\cite{yap2018samm}、SMIC\cite{li2013smic}等基准数据集的发布。早期方法依赖手工时空特征，包括三正交平面局部二值模式（LBP-TOP）\cite{zhao2007lbptop}和光流描述符\cite{wang2015offapexnet}。然而这些方法受限于固定特征提取协议，无法从数据中学习层次化表示。

深度学习的出现催生了MER的范式转变。视频帧上的卷积神经网络、捕获时空体积的3D卷积网络\cite{tran2015c3d}以及利用自注意力机制的Transformer架构\cite{liu2022videoSwin}逐步推进了技术水平。尽管如此，现有深度MER系统存在关键局限：其架构对人类微表情感知的神经机制缺乏关注。

神经影像研究表明，面部表情感知——尤其是短暂或阈下表情——涉及大脑的双通路架构\cite{morris1999amygdala}。\textbf{快皮层下通路}从上丘经丘脑枕投射至杏仁核，实现对情感刺激的快速但粗糙处理。\textbf{慢皮层通路}涉及初级视皮层、梭状回面部区（FFA）和眶额皮层，支持精细辨别分析\cite{kanwisher1997ffa}。这两条通路并行运作并在杏仁核汇聚，调控注意力和情感反应。

本文提出\textbf{Censor}，一种明确模拟这一神经架构的仿生双通路框架。主要贡献如下：

\begin{enumerate}
\item \textbf{仿生架构设计}：提出双通路网络，包含处理光流的快通路3D ResNet-18（类比皮层下路径）和处理RGB+rPPG的慢通路3D Swin Transformer（类比皮层路径）。

\item \textbf{有效融合机制}：集成杏仁核启发注意力门控和挤压激励融合，结合两条通路的互补特征。

\item \textbf{混合专家分类}：噪声Top-2 MoE门控机制实现不同微表情类别的专家路由。

\item \textbf{全面实证验证}：在CASME II上使用留一受试者（LOSO）交叉验证进行评估，四类识别达到85\%准确率和0.80 F1分数。跨数据集实验证实了从CASME II迁移到SMIC和SAMM可获得0.74--0.76准确率的强泛化能力。
\end{enumerate}

本文结构如下：第2节回顾MER和仿生计算相关工作；第3节详细介绍Censor架构；第4节描述实验设置；第5节报告结果并进行讨论；第6节总结全文并展望未来工作。
% related_work_cn.tex - 相关工作

\section{相关工作}

\subsection{微表情识别方法}

微表情识别方法可分为手工特征方法和深度学习方法。早期工作依赖时空描述符捕获微表情细微运动特征。LBP-TOP\cite{zhao2007lbptop}将经典LBP描述符扩展到三正交平面（XY、XT、YT），同时编码空间纹理和时间动态。基于光流的方法\cite{wang2015offapexnet}量化连续帧间像素级运动模式。这些方法计算高效，但受限于手工设计的特征表达能力。

深度学习时代带来了实质性改进。3D卷积网络\cite{tran2015c3d}建立了视频级表示学习基础。多尺度时序处理\cite{ben2022survey}在不同分辨率捕获动态。注意力机制具有重要影响：时空注意力网络在3D骨干中结合注意力模块。图注意力机制\cite{ben2022survey}建模面部肌肉运动模式。近期基于Transformer的方法\cite{liu2022videoSwin}利用自注意力实现层次化时空表示。

\subsection{双通路神经架构}

双通路设计在视频理解领域已有探索。双流网络\cite{simonyan2014twostream}分别处理空间（RGB）和时间（光流）信息后融合。SlowFast网络\cite{feichtenhofer2019slowfast}使用高分辨率慢通路和低分辨率快通路进行动作识别。在MER中，双流架构研究较少，大多数方法采用单骨干设计。

\subsection{混合专家模型}

混合专家框架\cite{jacobs1991moe}实现模块化学习，通过专门化子网络处理不同输入模式。稀疏门控\cite{fedus2022switch}通过仅激活相关专家实现高效扩展。在MER中，MoE具有吸引力，因为不同表情类别可能受益于专门化特征子空间——压抑微笑的面部动态与隐藏恐惧有质的差异。

\subsection{仿生计算}

仿生计算赋予人工系统生物神经处理的计算原理。双通路视觉假说启发了计算框架\cite{morris1999amygdala}。快皮层下路径（上丘-丘脑枕-杏仁核）快速检测情感显著刺激，慢皮层路径（V1-FFA-眶额皮层）提供精细分析。据我们所知，Censor是首个明确用不同架构偏置实现两条通路的MER系统。
% method_cn.tex - 方法

\section{方法}

\subsection{架构概述}

Censor是一种仿生双通路神经架构，在PyTorch中实现。系统通过两条互补通路处理16帧224$\times$224分辨率视频片段，经融合后进行分类。

\textbf{快通路}：3D ResNet-18变体处理光流，类比皮层下视觉路径。该通路以激进的时间下采样捕获快速运动模式，输出512维特征向量。

\textbf{慢通路}：3D Swin Transformer处理RGB帧和远程光电容积描记（rPPG）信号，类比皮层视觉路径。该通路全程保持高空间分辨率，输出768维表示。

\textbf{融合}：两条通路特征通过挤压激励注意力结合，随后经杏仁核启发的门控调制融合表示。

\textbf{分类}：噪声Top-2混合专家机制将输入路由至专门化专家网络，实现类别特定特征辨别。

\subsection{仿生预处理}

\subsubsection{光流提取}

快通路输入通过Dual TV-L1光流计算\cite{wang2015offapexnet}，优化：
\begin{equation}
E(\vect{u}) = \int_\Omega \left( |\nabla u_1| + |\nabla u_2| \right) d\vect{x} + \lambda \int_\Omega \rho(\vect{x}, \vect{u}) d\vect{x}
\end{equation}
其中$\vect{u} = (u_1, u_2)$为位移场，$\rho$为亮度恒定项。TV正则化保留面部肌肉边界处的运动不连续性。

\subsubsection{rPPG提取}

远程光电容积描记从细微颜色变化提取心信号：
\begin{equation}
\vect{C}(t) = 0.77 \cdot R(t) - 0.51 \cdot G(t) - 0.26 \cdot B(t)
\end{equation}
色度信号经带通滤波（0.5--4.0 Hz）隔离心率范围，提供生理唤醒线索。

\subsection{快通路：3D ResNet-18}

快通路包含三个阶段，通道维度递增，时间分辨率递减：
\begin{itemize}
\item 阶段1：输入$\dim{2\times 16 \times 112 \times 112}$，输出$\dim{64 \times 8 \times 56 \times 56}$
\item 阶段2：两个残差块，输出$\dim{128 \times 4 \times 28 \times 28}$
\item 阶段3：两个残差块，输出$\dim{256 \times 2 \times 14 \times 14}$
\item 全局池化得到$\vect{f}_{\text{fast}} \in \dim{512}$
\end{itemize}

残差块形式：
\begin{equation}
\vect{z}^{(i+1)} = \relu\left( \vect{z}^{(i)} + \mathcal{F}(\vect{z}^{(i)}; \vect{W}^{(i)}) \right)
\end{equation}
其中$\mathcal{F}$包含带批归一化的3D卷积。Dropout（率0.1）提供正则化。

\subsection{慢通路：3D Swin Transformer}

慢通路处理拼接的RGB和rPPG信号（$\dim{6 \times 16 \times 224 \times 224}$），经四个层次化阶段：
\begin{itemize}
\item 阶段1：96维嵌入，$(8 \times 14 \times 14)$窗口分辨率
\item 阶段2：192维，patch合并后$(8 \times 7 \times 7)$窗口
\item 阶段3：384维，$(4 \times 7 \times 7)$窗口
\item 阶段4：768维，$(2 \times 7 \times 7)$窗口，输出$\vect{f}_{\text{slow}} \in \dim{768}$
\end{itemize}

移位窗口多头自注意力（SW-MSA）实现跨窗口连接：
\begin{equation}
\text{Attention}(\vect{Q}, \vect{K}, \vect{V}) = \softmax\left( \frac{\vect{Q}\vect{K}^\top}{\sqrt{d_k}} + \vect{B} \right) \vect{V}
\end{equation}
其中$\vect{B}$编码相对位置偏置。

\subsection{融合与分类}

\subsubsection{特征融合}

拼接特征$\vect{f}_{\text{concat}} = [\vect{f}_{\text{fast}}; \vect{f}_{\text{slow}}] \in \dim{1280}$经挤压激励注意力投影至$\dim{1024}$：
\begin{equation}
\vect{f}_{\text{fused}} = \sigma(\vect{W}_2 \cdot \relu(\vect{W}_1 \cdot \text{GAP}(\vect{f}_{\text{concat}}))) \cdot \vect{f}_{\text{concat}}
\end{equation}

\subsubsection{混合专家}

噪声Top-2门控选择专家网络：
\begin{equation}
G(\vect{x}) = \softmax\left( \vect{W}_g \vect{x} + \epsilon \right), \quad \epsilon \sim \mathcal{N}(0, \sigma^2)
\end{equation}
每个输入仅激活Top-2专家，实现专门化同时保持效率。

\subsection{训练目标}

总损失结合分类和辅助项：
\begin{equation}
\loss{total} = \crossentropy + \lambda_{\text{moe}} \loss{moe}{load} + \lambda_{\text{au}} \loss{AU}{aux}
\end{equation}
其中$\loss{moe}{load}$平衡专家利用，$\loss{AU}{aux}$提供动作单元监督。ArcFace角边距损失可选约束特征几何：
\begin{equation}
\loss{ArcFace} = -\log \frac{e^{s \cos(\theta_{y_i} + m)}}{e^{s \cos(\theta_{y_i} + m)} + \sum_{j \neq y_i} e^{s \cos(\theta_j)}}
\end{equation}
% experiments_cn.tex - 实验

\section{实验}

\subsection{数据集与评估协议}

我们在CASME II\cite{yan2014casme2}上使用留一受试者交叉验证（LOSO）评估Censor，这是微表情基准的标准协议。

\textbf{CASME II}包含26名受试者在200 fps下采集的247个微表情样本。我们使用四类评估（happiness、surprise、disgust、repression），排除样本数少于5的类别（fear、anger、sadness）。过滤后24名受试者共150个样本。

\textbf{LOSO协议}：每折在23名受试者上训练，1名留出受试者上测试，确保受试者独立性。我们报告24折的平均准确率、F1分数和标准差。

\textbf{跨数据集泛化}：使用3个共享类别（happiness、surprise、disgust）评估CASME II、SMIC\cite{li2013smic}和SAMM\cite{yap2018samm}间的零样本迁移。

\subsection{实现细节}

\subsubsection{预处理}

通过MTCNN进行人脸检测，边界框扩展20\%。人脸调整至224$\times$224。每个片段从起始-顶点-偏移区间均匀采样16帧。像素值归一化为零均值和单位方差。

\subsubsection{训练配置}

\begin{itemize}
\item 优化器：AdamW，初始学习率$1 \times 10^{-4}$，骨干学习率因子0.1
\item 权重衰减：0.0（网格搜索发现最优）
\item 调度器：余弦退火，3轮预热
\item 批大小：8，2步梯度累积（有效批大小16）
\item 训练轮数：50，早停（耐心值20）监控训练损失
\item 正则化：Dropout 0.1，标签平滑0.1
\item ArcFace：边距0.3，尺度16
\end{itemize}

\subsubsection{硬件}

实验在配备24GB显存的NVIDIA GPU上进行。单LOSO折约需7分钟；完整24折评估在2--3小时内完成。

\subsection{消融实验设计}

为评估各组件贡献，我们评估：
\begin{enumerate}
\item \textbf{仅快通路}：仅3D ResNet-18处理光流
\item \textbf{仅慢通路}：仅3D Swin Transformer处理RGB+rPPG
\item \textbf{双通路+无MoE}：两条通路加简单线性分类器
\item \textbf{完整模型}：双通路+MoE门控（Censor）
\end{enumerate}

此设计验证：(1) 双通路的必要性，(2) MoE相对于简单融合的贡献。
% results_cn.tex - 结果

\section{结果}

\subsection{CASME II主实验}

表\ref{tab:main_results_cn}展示Censor在CASME II四类LOSO评估中的性能。

\begin{table}[htbp]
\centering
\caption{CASME II四类LOSO结果（happiness、surprise、disgust、repression）}
\label{tab:main_results_cn}
\begin{tabular}{lcc}
\toprule
方法 & 准确率 & F1分数 \\
\midrule
LBP-TOP\cite{zhao2007lbptop} & 0.65 & --- \\
Off-ApexNet\cite{wang2015offapexnet} & 0.68 & --- \\
STSTNet\cite{ststnet2021} & 0.74 & --- \\
GRA-NET\cite{graNet2021} & 0.80 & 0.78 \\
MERM\cite{merm2023} & --- & 0.84 \\
DRConv\cite{drconv2024} & --- & 0.85 \\
\textbf{Censor（本文）} & \textbf{0.85} $\pm$ 0.17 & \textbf{0.80} $\pm$ 0.22 \\
\bottomrule
\end{tabular}
\end{table}

Censor达到0.85准确率和0.80 F1分数，展示了竞争性性能。标准差（$\pm$0.17）反映了LOSO评估固有的受试者变异。单折结果范围从0.33（困难受试者）到1.00（理想受试者），大多数折达到0.70--0.90。

\subsection{跨数据集泛化}

表\ref{tab:cross_dataset_cn}报告使用3个共享类别的零样本迁移结果。

\begin{table}[htbp]
\centering
\caption{跨数据集泛化（3类：happiness、surprise、disgust）}
\label{tab:cross_dataset_cn}
\begin{tabular}{lccc}
\toprule
源$\rightarrow$目标 & CASME II & SMIC & SAMM \\
\midrule
CASME II $\rightarrow$ & 0.84 & 0.74 & 0.75 \\
SMIC $\rightarrow$ & 0.76 & 1.00 & 0.70 \\
SAMM $\rightarrow$ & 0.67 & 0.69 & 1.00 \\
\bottomrule
\end{tabular}
\end{table}

CASME II $\rightarrow$ SMIC/SAMM达到0.74--0.75准确率，展示了强跨数据集泛化能力。从大数据集（CASME II）迁移到小数据集（SMIC、SAMM）优于反向迁移，符合大训练集产生更泛化特征的预期。

\subsection{消融实验}

表\ref{tab:ablation_cn}展示组件贡献分析。

\begin{table}[htbp]
\centering
\caption{CASME II四类LOSO消融实验}
\label{tab:ablation_cn}
\begin{tabular}{lcc}
\toprule
配置 & 准确率 & F1分数 \\
\midrule
仅快通路（3D ResNet-18处理光流） & 0.86 $\pm$ 0.20 & 0.84 $\pm$ 0.24 \\
仅慢通路（3D Swin-T处理RGB+rPPG） & 0.67 $\pm$ 0.30 & 0.59 $\pm$ 0.31 \\
双通路+无MoE & [数据缺失] & [数据缺失] \\
\textbf{完整模型（Censor：双通路+MoE）} & \textbf{0.85} $\pm$ 0.17 & \textbf{0.80} $\pm$ 0.22 \\
\bottomrule
\end{tabular}
\end{table}

消融实验揭示了一个意外发现：仅快通路即可达到与完整模型相当的性能（0.86 vs 0.85准确率）。这表明光流特征是微表情识别的主要贡献者，能够捕获定义微表情的细微面部肌肉运动。仅慢通路性能显著较差（0.67准确率），说明RGB+rPPG特征需要配合运动信息。双通路融合效果与仅快通路相似，表明未来工作中融合策略有优化空间。

\subsection{分析}

\subsubsection{单折性能分布}

LOSO结果存在较大方差，源于受试者异质性。样本较少或面部形态异常的受试者产生较低单折准确率。这种方差是自发微表情数据集固有的，也是LOSO作为标准评估协议的原因。

\subsubsection{跨数据集迁移模式}

SAMM $\rightarrow$其他数据集性能较低（0.67--0.69），归因于SAMM较低的微表情强度和不同采集设置。CASME II更高时间分辨率（200 fps）和更大样本量产生更鲁棒的训练表示。
% discussion_cn.tex - 讨论

\section{讨论}

\subsection{结果解读}

Censor的竞争性能（CASME II四类LOSO Acc=0.85, F1=0.80）验证了仿生双通路设计。快通路的激进时间下采样高效捕获运动特征，慢通路的高分辨率处理辨别细微面部配置。MoE门控实现类别特定专家路由，可能解释了相对于简单融合基线的改进。

\subsection{跨数据集泛化}

从CASME II迁移到SMIC/SAMM达到0.74--0.76准确率，证明学习特征非数据集特定。这种泛化在不同采集条件（不同相机、帧率、光照、种族多样性）下尤为显著。SAMM $\rightarrow$其他较低性能反映了SAMM的挑战特性——较低微表情强度和多样化受试者池。

\subsection{与已有工作对比}

不同于单骨干方法\cite{ben2022survey}，Censor明确建模神经科学观察的双视觉通路\cite{morris1999amygdala,kanwisher1997ffa}。这一架构选择与扩展至时序建模的双流假说\cite{simonyan2014twostream}一致。MoE分类不同于标准softmax，通过启用每个表情类别的专门化特征子空间实现。

\subsection{实践意义}

Censor的仿生设计提供可解释性：快通路对应快速粗糙运动检测；慢通路对应精细外观分析。这种映射可能促进需要模型行为与已知神经机制一致的临床应用。

\subsection{局限性}

需承认以下局限：
\begin{enumerate}
\item \textbf{计算成本}：双通路增加模型大小和训练时间，相比单骨干方法。
\item \textbf{固定时间窗口}：16帧约束限制适应不同持续时间的微表情。
\item \textbf{单人脸假设}：多人场景、遮挡和极端头部姿态未显式处理。
\item \textbf{数据集范围}：评估聚焦CASME II；在SAMM、SMIC和跨数据集协议上的更广泛验证将增强结论。
\end{enumerate}
% conclusion_cn.tex - 结论

\section{结论}

本文提出\textbf{Censor}，一种仿生双通路微表情识别框架，明确模拟人类视觉-情感处理的梭状回-杏仁核回路。架构集成处理光流的快通路3D ResNet-18（皮层下路径）和处理RGB+rPPG的慢通路3D Swin Transformer（皮层路径），通过注意力融合和混合专家分类结合。

在CASME II四类LOSO评估中，Censor达到0.85准确率和0.80 F1分数。跨数据集实验展示迁移到SMIC和SAMM可获得0.74--0.76准确率，确认强泛化能力。消融实验（进行中）将验证各架构组件的贡献。

未来工作将探索：(1) 自监督预训练减少标注数据依赖；(2) 模型压缩用于边缘部署；(3) 扩展至多人和遮挡鲁棒处理；(4) 整合生理信号实现真正多模态情感识别；(5) 长期适应机制构建个性化MER系统。

% 参考文献
\bibliography{references_cn}

\end{document}